\section{A finite field request before any choice of column}

% BEGIN F77W1
\paragraph{F77W1. Residue roots for all selected twists.}
Fix the prehost and the allowed character set $\Xi$ at the second
field request. For each $\xi\in\Xi$, each finite prime $Q$ of $E$
where $\xi$ ramifies, and each $P$ dividing the nonzero integer
$NQ$, request a primitive root of unity of order $\#k(P)$.
Use $\exp(2\pi i/\#k(P))$. This is a finite set of algebraic
numbers for finite $\Xi$: K77F1 gives finiteness of each ramification
set, and the ideal generated by $NQ$ has finitely many prime factors.
The request is defined as empty on infinite $\Xi$, since no admitted
column's finite character range can equal such a set.

To check coverage take a later column $I$ with character range
$\Xi$ and $P\in I.Q(a)$. The place $v$ of the real field below
$P$ is one at which the finite tuple character is nontrivial on
the subgroup of integral units supported over $v$. Some prime
$Q$ over $v$ ramifies for the character $\xi_a$ which reads that
tuple. Otherwise every local factor on the indicated subgroup
would be one. Indeed the components at primes above $v$ and their
inverses are integral, and at the other primes the component is
one. The finite idele product formula and the character reading
would make the entire subgroup trivial. Now $NQ\in Q$ and
contraction to $v$ followed by extension into $P$ gives $NQ\in P$.
The requested set therefore contains the primitive root of order
$\#k(P)$. Membership in the distinguished image of $L$ pulls it
back to a primitive root in $L$ by injectivity. The other field
clause of $\mathrm{ValsIn}$, the unramified values of $\xi_a$,
is exactly the column's character reading on uniformizers. This
proves $\mathrm{ValsIn}$ simultaneously for every admitted column.
% END F77W1

% BEGIN F77W2
\paragraph{F77W2. Splitting the row and lift operators before the column.}
Keep the prehost, the frame, the two fields in the request and
$\Xi$ fixed. The tuple of configuration data used for the Hecke
operators has three entries: the two marked generators $\pi_g$,
the residue character function and the selected prime sets by
character. Each $\pi_g$ is an algebraic integer of the fixed host
with $|\tau(\pi_g)|^2=g$ for every embedding $\tau$. The
integral lattice in archimedean space has finitely many points
with these bounds. The residue character on each marked prime
is fixed by the inverse of the base character on unit residues
(and is zero in the unused cases of the stored function).
For each $\xi\in\Xi$ the selected primes are among the primes
above the finitely many real places where the finite character
is nontrivial on the supported units, by the argument of F77W1.
There are finitely many subsets, independently at the finitely
many occupied indices, and the other entries are empty. Thus
finitely many triples occur. Retain one admitted representative
of each realized triple.

At a representative $I_0$ the ordinary row and each lift row
have a finite-dimensional equivariant function space over its
coefficient field $L_0'$. For each of their Hecke families choose
a finite set of operators from the family spanning its linear
span in the endomorphism space. Take the product of their minimal
polynomials, or one for an empty set. It annihilates the chosen
generators. Request the images of a rational basis of $L_0'$ and
the complex roots of these polynomials for every tuple and required
row, unioned over the representatives. All these requests are
finite and algebraic.

For a new admitted column, inclusion of the basis gives the unique
field inclusion $L_0'\to L$ compatible with the two distinguished
embeddings. Its tuple matching a representative is found by equality
of the actual complex Hecke characters; equality of configuration
gives the identical marked generators and prime sets. Evaluating
the character reading after field inclusion identifies the tuple
characters, and evaluating on lifts of residue units identifies the
residue characters. Hence the levels, their determinant characters,
and the glued ordinary and lift types are scalar extensions after
reindexing the same prime set. On the equivariant functions this
base change can be checked on coordinates in an $L_0'$-basis of
$L$: expanding a function gives equivariant coefficient functions,
and the converse follows by linearity. Double-coset sums commute
with the inclusion coefficient by coefficient. The annihilating
polynomials for the chosen generators split in $L$ by the root
request. For commuting operators this supplies an invariant flag
over $L$: take successive kernels of split linear factors to find
a common eigenline, then pass to the quotient and induct on
dimension. Linear combinations of the generators preserve that
flag. It follows that every member of the family has split
characteristic polynomial, with the zero-dimensional case given
by the empty flag. To obtain exactly $\mathrm{SplitOver}$, write
the characteristic polynomial of any one such operator $T$ as
$\prod_\lambda (z-\lambda)^{m_\lambda}$ over $L$. Cayley--Hamilton
annihilates $T$ by this product. The factors for distinct eigenvalues
are coprime, so their B\'ezout identities decompose the whole space
as the direct sum of $\ker(T-\lambda)^{m_\lambda}$. In particular
the maximal generalized eigenspaces span the whole space for every
operator, including one with repeated eigenvalues. This uses the
original commutation laws for the ordinary and lift Hecke families,
including their inert operators.

Taking the union of this set, F77W1 and U76W2 at the second request
gives the original character values, the splitting on every tuple
and the Weber fields needed in Q76Z1. The distinguished complex
place respects conjugation by the base involution. In conjunction
with the absolute irreducibility of the selected character set and
the above commutation laws, these are the hypotheses of the
three-stage family construction $\mathrm{fam3.Ok3}(1)$.
The tuple is always one of the column's original tuples, including
the base tuple, so both row and lift assertions apply to it.
The union is taken before the prime selection, column, trace, group
characters and the later analytic records.
% END F77W2
